\graphicspath{{chapters/04design/graphs}}

\chapter{Design and Implementation}

%This chapter may be called something else\ldots but in general the
%idea is that you have one (or a few) ``meat'' chapters which describe
%the work you did in technical detail.

This chapter outlines the steps taken to design and implement a new scheduler for LLM workloads, with support for Mixture-of-Experts models on heterogenous hardware.
We begin by outlining the problem description of an LLM scheduler, introducing the terminology used throughout this section. 

Next, we explain the choice of simulator used for modelling the system's expected throughput for a given workload.
We then use this simulator to define the objectives for a bounded-constraint optimisation problem, which represents a subset of the overall scheduler.
Then, we describe several scheduling assumptions that can be used as optimisation strategies in the complete scheduler.

We finish by bringing together the workload modelling, objectives, constraint optimisation and scheduling assumptions into a single problem formulation.
Here, we separate the scheduler into a two-stage problem, with an \textit{inner} and \textit{outer} loop.
The \textit{inner loop} assigns specific workloads to GPU islands using deterministic bin-packing.
In contrast the \textit{outer loop} performs a heuristic search to find how best to slit up the GPU inventory into islands.
Using this formulation, we describe the specific design and then implementation of the two loops, in preparation for evaluation and testing.

\section{Problem Description}

The purpose of an LLM scheduler within an inference system, is to organise and configure a collection of GPUs for optimal performance.
We can utilise previous \textit{phase splitting} technology, to assign the prefill and decode phases of inference to optimal hardware, given the significant efficiency gains. 

To create a schedule we must first subdivide a list of GPUs into \textit{islands}, an island is a collection of homogenous GPUs of a given type and size.
Presented with a list of islands, we must then assign either prefill or decode work (and configure appropriate parallelism levels with respect to MoE) to these islands, maximising a defined objective.
In this case we are optimising for cost efficiency, so throughput with respect to cost is a good metric.

The logical architecture of such inference system is displayed in \autoref{fig:04-service-architecture}, where incoming requests are mapped to prefill islands.
Then, after transferring the (KV) cache, completed prefill requests are sent to the optimum decode island.
We are developing a new scheduler component for this system.

\begin{figure}[ht]
    \centering
    \begin{tikzpicture}[node distance=1cm, auto]
        \node[block] (rt) {Traffic Router};
        \node[block, right=1cm of rt, yshift=2cm, text width = 2cm] (p1) {Prefill\\Island};
        \node[block, right=1cm of rt, yshift=-0.5cm, text width = 2cm] (p2) {Prefill\\Island};
        \node[block, right=1cm of p1, yshift=-1cm, text width = 2cm] (d1) {Decode\\Island};
        \node[block, right=1cm of p2, yshift=-1cm, text width = 2cm] (d2) {Decode\\Island};
        
        \node[block, left=1cm of rt, yshift=2.75cm] (sec) {Scheduler};
        
        % landing spots on west face of router
        \coordinate (in1) at ($(rt.west)+(0cm,0.3cm)$);
        \coordinate (in2) at ($(rt.west)+(0cm,0.15cm)$);
        \coordinate (in3) at ($(rt.west)+(0cm,-0.15cm)$);
        \coordinate (in4) at ($(rt.west)+(0cm,-0.3cm)$);

        % spaced-out input arrows
        \draw[->, thick] ([xshift=-1cm,yshift=1cm] rt.west)  -- (in1);
        \draw[->, thick] ([xshift=-1cm,yshift=0.5cm] rt.west) -- (in2);
        \draw[->, thick] ([xshift=-1cm,yshift=-0.5cm] rt.west)-- (in3);
        \draw[->, thick] ([xshift=-1cm,yshift=-1cm] rt.west) -- (in4);
        
        % label for those four arrows
        \node[anchor=east] at ($(rt.west)+(-1.2cm,0cm)$) {User Requests};
        
         % dotted right-angle arrow from Scheduler to Traffic Router
\draw[->, thick, dotted]
    ([yshift=-0.15cm] sec.east) -| (rt.north);
        
        % Fit box around all islands
        \node[draw, dotted, inner sep=0.6cm, fit=(p1)(p2)(d1)(d2)] (Islands) {};

        % compute the shifted start/end of the dotted arrow
        \coordinate (arrStart) at ($(sec.east)+(0,0.15cm)$);
        \coordinate (arrEnd)   at ($(sec.east -| Islands.west)+(0,0.15cm)$);

        % single straight, horizontal dotted arrow from Scheduler into box (offset up)
        \draw[->, thick, dotted] (arrStart) -- (arrEnd);
        
        % direct arrows from router to prefills
        \draw[->, thick] (rt.east) -- (p1.west);
        \draw[->, thick] (rt.east) -- (p2.west);

        % direct arrows from prefills to decode islands
        \draw[->, thick] (p1.east) -- (d1.west);
        \draw[->, thick] (p1.east) -- (d2.west);
        \draw[->, thick] (p2.east) -- (d1.west);
        \draw[->, thick] (p2.east) -- (d2.west);

    \end{tikzpicture}
    \caption[Sample service architecture of LLM inference engine]
    {Sample service architecture of LLM inference engine, the scheduler components directs how traffic is distributed across the GPUs and confgiures the islands}
    \label{fig:04-service-architecture}
\end{figure}

\let\clearpage\relax
\section{Simulator}

In order to be able to develop an LLM scheduler, we need access to many high-performance GPUs to test our schedule configurations.
Unfortunately, GPU-hours are prohibitively expensive and using such hours for ``dummy'' workloads, to test a scheduler, would be considered a waste of resources. %CITE
Instead, we utilise the latest simulator technology to estimate the performance of LLM inference given a set of input parameters.
Several simulators already exist that attempt to model the performance of an LLM request by using estimation techniques.

% Previous work has also used schedulers too (reference the work here).

\textbf{Vidur} \autocite{agrawalVidurLargeScaleSimulation2024} is an LLM simulator developed by Microsoft, designed to address the high cost and complexity of experimentally finding the best deployment configuration of LLMs.
This is achieved by profiling a model-gpu combination for expected throughput on given operations (different prefill and decode requests).
Then, using a small machine learning model, the system predicts the likely performance with high degrees of accuracy for different parrallelism configurations.

\textbf{Shallowsim} \autocite{icezack12Icezack12Shallowsim2025} is a specialist simulator designed to estimate the compute times for different LLM request types.
Specifically built for DeepSeek-V3/R1 models, it has native support for simulating MoE architectures.
Unlike Vidur, Shallowsim requires no profiling phase, rather a list of device performance metrics is provided e.g. \texttt{fp8} performance.
These metrics are then used to compute the expected performance for the prefill and decode phase separately

We select Shallowsim for its transparency of calculations (simplicity), as it uses no machine learning to predict performance.
Instead, simple mathematical operations are used, which allows us to inspect performance and modify the calculations as necessary.
The Shallowsim simulator also already has native support for Mixture of Expert models, which is our key novel design goal for the scheduler.

%And requires complex and expensive profiling

\subsection{Shallowsim Implementation}

%\todo{
%\textit{This is just describing how shallowsim does its maths, will finish very last as it is not strictly necessary for the report - Might even include in the appendix if too long.}
%}

%\textbf{Prefill}
%
%Dense-MLA (dense Multi-head/Multi-query Attention compute),
%TP-MLA (Dense-MLA plus the Tensor-Parallel all-reduce),
%Dense-MLP (feed-forward network compute),
%Shared-Expert (MoE expert run that every token visits),
%Routed-Expert (MoE expert run that only selected tokens visit),
%Every token first passes through the Shared Expert, but a subset of tokens are also routed to extra specialists. This is modelled with the Routed Expert function.
%%Dispatch (All-to-All that ships tokens out to remote experts),
%%Combine (All-to-All that brings tokens back),
%%Sending activations to the respective routed expert nodes is modelled using the Dispatch step, the Combine step then models how processed activations are collated back together. These two functions model the network overhead of LLM queries.
%
%\textit{Compute Time
%-	Dense MLA, Dense MLP (Dense Layers)
%-	TP-MLA, Shared Expert, Routed Expert (Sparse Layers)
%Communication Timemake 
%-	Dispatch, Combine (Sparse Layers)
%Overlap Penalties
%-	If Combine is longer than TP-MLA + Shared Expert
%-	If Dispatch is longer than Routed Expert
%Sum
%-	Compute + Communication + Any Overlap = Prefill Time
%}
%
%\subsubsection{Decode}
%
%Load-KV,
%Dense-MLA (as prefill),
%Sparse-MLA,
%Dense-MLP (as prefill),
%Shared-Expert (as prefill),
%Routed-Expert (as prefill),
%Dispatch (as prefill),
%Combine (as prefill),
%
%\textit{Compute Time
%-	Dense MLA, Dense MLP (Dense Layers)
%-	Sparse MLA, Shared Expert, Routed Expert (Sparse Layers)
%Communication Time
%-	Dispatch, Combine (Sparse Layers)
%Overlap Penalties
%-	TODO
%Sum
%-	TODO = Per token decode time 
%}

\subsection{Modifactions}

We make several small modifications to the Shallowsim simulator to support our use-case of efficiently profiling system configurations.
As-built Shallowsim is designed to profile several different parallelism configurations and batch sizes, then report back the overall fastest configuration.
We need support for profiling a specific parallelism configuration and batch size, as we may not always be able to use the most optimal for a given exact sequence length, but rather a range of sequence lengths.

To do this, we extract all the mathematical operations present in the existing code and form procedures with specific input parameters, e.g., $\texttt{prefill\_a2a}(\texttt{gpu}, \texttt{seq\_len}, \texttt{tp\_num})$.
Here, we precisely specify the gpu model, input sequence length, and tensor parallelism degree.

Given the modifications, shallowsim now supports the following functions:
\begin{itemize}
    \item \textbf{Prefill Performance}\ \texttt{-> Milliseconds}\\
    \texttt{prefill\_time(gpu, seq\_len, kv\_cache\_rate, tp\_num, dp\_num)}
    
    \item \textbf{Decode Performance}\ \texttt{-> Milliseconds}\\
    \texttt{decode\_time(gpu, tp\_num, bs\_num, seq\_len, decode\_len, moe\_num, device\_num)}
    
    \item \textbf{Maximum Batch Size}\  \texttt{-> Int}\\
    \texttt{check\_max\_bs(gpu, tp\_num, bs\_num, seq\_len, decode\_len, moe\_num, device\_num)}
\end{itemize}

\subsection{Usage}
\label{chap:simulator-usage}

\subsubsection{Prefill Phase}
Using the Shallowsim simulator we can measure the time $t$ (in milliseconds) to compute a single prefill request for a given set of parameters, and convert this to throughput R (in requests per second) via
$R_{\text{prefill}} = \frac{1000}{t}$.

% include sample parameters for prefill function

\subsubsection{Decode Phase}
Shallowsim has support for calculating the time in milliseconds for a single token generation step.
As we would typically expect several tokens to be generated, we cannot directly compute requests per second using Shallowsim, instead we must perform additional computations.
The decode time of a single token is affected by the input/current sequence length, therefore we must measure the single token compute time for every token to be generated in a request.

% graph of sequence length affecting decode time

By measuring $t_i$ the compute time (in milliseconds) for the $ith$ token, the total decode time for a request generating N tokens is $t_{\text{decode}} \;=\;\sum_{i=1}^{N} t_i\,$.
We can then convert this total time into throughput $R_{\text{decode}}$ (in requests per second), using the same method as prefill.
$R_{\text{decode}}
\;=\;\frac{1000}{t_{\text{decode}}}
\;=\;\frac{1000}{\sum_{i=1}^{N} t_i}\,$,
where $R_{\text{decode}}$ represents how many full decode requests per second the configuration can sustain.


% How requests will be modeled
\section{Workload Modeling}
\label{chap:workload-modeling}

Statistical distributions (available from sources)

\subsection{Input Sequence Length}
To achieve more fine-grained control of request scheduling, dependant on input sequence, we model the input request length as a probability distribution.
Hence, for a chosen workload type, there is a specific probability of request arriving with given sequence length.

Let $w \in W$ denote a workload type from the set of all workload types $W$.
For each workload type $w$, we define a probability mass function $p(s \mid w)$, which represents the probability of a request having sequence length $s \in S$ under workload $w$.
This distribution satisfies:

\[
\sum_{s \in S} p(s \mid w) = 1 \quad \text{for all } w \in W
\]

% Equation of PDF

% Graph of PDF
%\begin{figure}[ht]
%    \centering
%    \includegraphics[width=\textwidth]{sample.pdf}
%    \caption[Sample probability distribution for input sequence lengtht]
%    {Sample probability distribution for input sequence length}
%    \label{fig:04-sample-pdf}
%\end{figure}

Instead of only considering the distribution of sequence lengths during assignment, to model system throughput, we use them to find the most optimal pairing of sequence length to GPU island.
Therefore, specific GPU islands can be configured to most optimally perform prefill operations for a specific sequence length.
When a request arrives, we can direct it to the GPU islands targeted to serve that request length, defined in the schedule.

% Graph of sequence length vs GPU
%\begin{figure}[ht]
%    \centering
%    \includegraphics[width=\textwidth]{sample.pdf}
%    \caption[Stuff]
%    {Stuff}
%    \label{fig:04-sequence-length-gpu}
%\end{figure}

\subsection{Decode Length}
While decode length can also be modelled using a probability density function, this information (specific distribution) is less useful from a scheduling perspective.
In contrast to input sequence length, the decode length is not known the moment a request arrives.
Instead, the generation process proceeds token-by-token until an explicit or implicit stop condition is met (e.g. end-of-sequence token). 

We can assume that a decode request must complete on the GPU island it started processing on, as switching island would incur additional communication overhead and significantly impact the request response time. %Why?

For some workloads there is a small correlation between input sequence length and decode length, but this relationship is often not significant. %Example?
Therefore, at the point of request arrival we do not have enough information about the decode length to make informed per-decode length scheduling decisions.

Hence, creating specific GPU configurations for short, medium, long, etc. decode lengths would likely be inefficient. 
As we cannot accurately predict decode length, there is a high probability that a decode request would have to move island.

Prior work has demonstrated that it can be possible to predict decode length using the input text contents. %Reference?

Instead, we model our scheduler using an average decode length for a given workload type and design every GPU decode island to support the maximum decode length expected. 
Unlike the sequence length, the decode length is simply being used for system modelling activities to characterise performance, not direct traffic.

However, we can still make scheduling decisions for the decode phase using the sequence length as a feature, given the maximum decode length affected by the input length.
This is due to all the input and decode tokens KV values being stored in GPU memory.

% Objectives
\section{Objectives}

%We will use requests per second as our measure of throughput for the overall system as it can be used as an overall measure of performance
%
%encouraging a cost effective solution


\subsection{Throughput Modelling}

Demonstrated in \autoref{chap:simulator-usage}, the Shallowsim simulator has support for estimating the throughput of a single request of a static input sequence and decode length.
It is highly improbable that an entire workload consists of a single request type, instead we expect a distribution of sequence lengths, as shown in \autoref{chap:workload-modeling}.
Hence, we need a method to accurately calculate the throughput of a system, with a variety of input sequence.

Using the probability distribution of the workload, we can compute the weighted average of throughput, using measured performance values for each sequence length.
Let $R_{\text{prefill}}(s)$ denote the measured prefill throughput (in requests per second) for a request of sequence length $s$.
The expected prefill throughput for workload $w$, denoted $\mathbb{E}[R_{\text{prefill}} \mid w]$ (expectation), is then given by:

\[
\mathbb{E}[R_{\text{prefill}} \mid w] = \sum_{s \in S} R_{\text{prefill}}(s) \cdot p(s \mid w)
\]

In addition to prefill, we can also compute the expected throughput for the decode phase. 
While prefill throughput depends solely on the input sequence length $s$, decode throughput typically depends on both the input length $s$ and the output (decode) length $d$.
For simplicity, we assume a fixed average decode length $d$ as a parameter.
Let $R_{\text{decode}}(s, d)$ denote the measured decode throughput (in requests per second) for an input sequence length $s$ and average decode length $d$.
Given a workload-specific sequence length distribution $p(s \mid w)$, the expected decode throughput for workload $w$ is:

\[
\mathbb{E}[R_{\text{decode}} \mid w, d] = \sum_{s \in S} R_{\text{decode}}(s, d) \cdot p(s \mid w)
\]

To account for the complete request processing pipeline, we define the overall expected throughput as the minimum of the expected prefill and decode throughputs.
Since requests cannot be served faster than the slowest stage, the effective throughput is bounded by the bottleneck.
Let $\mathbb{E}[R_{\text{prefill}} \mid w]$ denote the expected prefill throughput for workload $w$, and $\mathbb{E}[R_{\text{decode}} \mid w, d]$ the expected decode throughput for workload $w$ and average decode length $d$.
The overall expected throughput is given by:

\[
\mathbb{E}[R_{\text{overall}} \mid w, d] = \min\left( \mathbb{E}[R_{\text{prefill}} \mid w],\ \mathbb{E}[R_{\text{decode}} \mid w, d] \right)
\]

This value represents the effective throughput achievable by the system under the given workload and output length assumptions.

\begin{figure}[ht]
    \centering
    \includegraphics[width=\textwidth,height=5cm]{example-image-a}
    \caption[Sample expected throughput]
    {Sample expected throughput, calculated using the workload shown in \autoref{fig:04-sample-prefill-pdf}.}
    \label{fig:04-sample-expected-throughput}
\end{figure}

\subsection{Throughput Allocation}

Previous calculations assumed that a singular GPU island is contributing to the overall throughput of the system.
In reality, we would like to be able to construct several GPU islands and assign each island to a request type it is best suited for.
Hence, we allocate a percentage of GPU resources for a given island to a given sequence length it specialises in.
Let \(I\) be the set of all GPU islands in the system, which we partition into two separate subsets \(I_{\mathrm{prefill}}\) and \(I_{\mathrm{decode}}\).
Let \(x_i(s)\) denote the fraction of overall GPU capacity on island \(i\) devoted to handling requests of length \(s\). 
We normalise these allocations so that \(\sum_{s\in S} x_i(s) = 1 \quad\text{for every }i\in I. \)
By including \(x_i(s)\) into our previous per-length throughput models, the expected prefill throughput under workload \(w\) becomes
\[
\mathbb{E}[R_{\mathrm{prefill}}\mid w]
=\sum_{i\in I_{\mathrm{prefill}}}\sum_{s\in S}
x_i(s)\,R_{\mathrm{prefill}}^{(i)}(s)\,p(s\mid w),
\]
while the expected decode throughput for average decode length \(d\) is
\[
\mathbb{E}[R_{\mathrm{decode}}\mid w,d]
=\sum_{i\in I_{\mathrm{decode}}}\sum_{s\in S}
x_i(s)\,R_{\mathrm{decode}}^{(i)}(s,d)\,p(s\mid w).
\]

\begin{figure}[ht]
    \centering
    \includegraphics[width=\textwidth,height=5cm]{example-image-a}
    \caption[Sample throughput allocation mapping]
    {Sample throughput allocation mapping.}
    \label{fig:04-sample-throughput-allocation}
\end{figure}

% Some graphs showing the throughput calculation + Some graphs showing throughput allocation outcomes

% Constraints
\section{Constraint Optimisation}
\label{chap:constraint-optimisation}

We can combine the previous throughput calculations to form a simple constraint optimisation problem with an objective to maximise the total expected throughput $R$ and minimise the deviation $\Delta$ from the target workload distribution $p(s\mid w)$.

\subsection{Problem Formulation}

The following statement describes a sample non-linear optimisation problem that can be solved for either prefill or decode throughput.

\subsubsection*{Input Data}
\[
\begin{aligned}
& I &&\text{(set of GPU islands)},\\
& S &&\text{(set of sequence lengths)},\\
& b_i(s) &&\text{(benchmark throughput of island \(i\) on length \(s\))},\\
& p(s\mid w) &&\text{(target probability of length \(s\) under workload \(w\)),}
\quad \sum_{s\in S}p(s\mid w)=1,\\
& \lambda &&\text{(trade-off parameter, $\lambda>0$, balancing throughput vs.\ deviation)}.
\end{aligned}
\]

\subsubsection*{Decision Variables}
\[
x_i(s)\;\ge\;0
\quad\text{(fraction of island \(i\)'s GPU capacity devoted to length \(s\))}.
\]

\subsubsection*{Constraints}
\[
\sum_{s\in S}x_i(s)=1
\quad\forall\,i\in I.
\]

\subsubsection*{Intermediate Quantities}
\begin{enumerate}
  \item \textbf{Performance per cell:}
  \[
    r_i(s) = b_i(s)\,x_i(s).
  \]
  \item \textbf{Total performance:}
  \[
    R = \sum_{i\in I}\sum_{s\in S}r_i(s)
      = \sum_{i\in I}\sum_{s\in S}b_i(s)\,x_i(s).
  \]
  \item \textbf{Performance by sequence length:}
  \[
    R_s = \sum_{i\in I}r_i(s)
        = \sum_{i\in I}b_i(s)\,x_i(s),
    \quad\forall\,s\in S.
  \]
  \item \textbf{Achieved share:}
  \[
    \sigma_s = \frac{R_s}{R}.
  \]
  \item \textbf{Deviation from target:}
  \[
    \delta_s = \bigl|\sigma_s - p(s\mid w)\bigr|,
    \qquad
    \Delta = \sum_{s\in S}\delta_s.
  \]
\end{enumerate}

\subsubsection*{Objective}
Maximise
\[
R \;-\;\lambda\,\Delta,
\]
where
\[
R = \sum_{i\in I}\sum_{s\in S}b_i(s)\,x_i(s),
\quad
\Delta = \sum_{s\in S}\bigl|\sigma_s - p(s\mid w)\bigr|,
\quad
\sigma_s = \frac{\sum_{i\in I}b_i(s)\,x_i(s)}{R}.
\]

\subsection{Applicable Solvers}

Non-linear (due to multiplication and division).

Can be solved using Python but is slow. 

\subsection{Problem Usage}

Cannot be used in isolation, we need an architecture that is able to formulate islands.

Turn unbounded bin packing to bounded bin packing.

% Heuristics used
\section{Scheduling Assumptions}

We can make several assertions and assumptions in an attempt to reduce the search space of the scheduler generation process.
These effectively apply heuristics to our problem, that are crafted to simplify the computation requirements while retaining solution quality.

\subsection{Prefill Parallelism}
\label{chap:prefill-parallelism}

To reduce the search space of the scheduler, we can assume for a given model and GPU type, there exists an optimal parallelism configuration that is consistent across input sequence lengths.
Demonstrated in \autoref{fig:04-prefill-throughput-vs-sequence-length} we can see that for all GPUs, $tp=64$ and $dp=1$ is the fastest parallelism configuration for all sequence lengths above 128 tokens.
Given we expect the majority of input requests to be over 128 tokens, we can reason that there is an \textit{optimal} GPU prefill configuration.

\begin{figure}[ht]
    \centering
    \includegraphics[width=\textwidth]{simulator_prefill_throughput_vs_seq_len.pdf}
    \caption[Stuff]
    {Stuff}
    \label{fig:04-prefill-throughput-vs-sequence-length}
\end{figure}

% Add some explanations of the graph

By pre-computing this optimal configuration, we significantly reduce the number of parameters considered by the scheduler.
This is because an island cannot be formed with several different parallelism configurations simultaneously. %Why?
And hence, the scheduler would have to contain a variable for each sequence-length-parallelism combination.

To identify a best GPU prefill configuration across sequence lengths we use the following calculations, let \(G\) be the total number of GPUs and consider the set of all parallelism configurations $C = \{\,c = (\mathit{tp},\mathit{dp}) \mid \mathit{tp}\cdot\mathit{dp} = G\}$.
For each configuration \(c\in C\) and sequence length \(s\in S\), we measure the prefill throughput \(R_{\mathrm{prefill}}(s,c)\) and define the optimal throughput at length \(s\) by
\[
R_{\mathrm{prefill}}^*(s) = \max_{c\in C} R_{\mathrm{prefill}}(s,c).
\]
We then compute the relative deviation of configuration \(c\) at length \(s\) as
\[
\delta(c,s) = \frac{R_{\mathrm{prefill}}^*(s) - R_{\mathrm{prefill}}(s,c)}{R_{\mathrm{prefill}}^*(s)}.
\]
Aggregating these deviations over all lengths gives
\[
\begin{aligned}
\bar{\delta}(c) &= \frac{1}{\lvert S\rvert} \sum_{s\in S} \delta(c,s), 
&&\text{(uniform average over \(S\))},\\
\widetilde{\delta}(c) &= \sum_{s\in S} \delta(c,s)\,p(s\mid w),
&&\text{(weighted by workload distribution)}.
\end{aligned}
\]
Finally, the best overall configuration is chosen as the one minimising the average deviation from the individual best configurations.
\[
c_{\mathrm{base}} = \arg\min_{c\in C} \bar{\delta}(c).
\]

\subsection{Decode Parallelism}
\label{chap:decode-parallelism}

We can compute the same best configuration for the decode phase, with small modifications to the set of $C$ configurations taking into account mixture-of-experts.
Unlike prefill, \(dp\) (the data‐parallelism factor) is not a native parameter of the Shallowsim simulator.
Instead, we precompute how many GPUs to pass to the simulator and then scale the resultant throughput by \(dp\).
Let \(G\) be the total number of GPUs, \(tp\) tensor parallelism factor, and \(n_r\) the total number of routed experts in the model.
We can form every pairing of valid factors using:
\[
T = \{\,tp \in \{1,\dots,G\}\mid G \bmod tp = 0\}, 
\qquad
D = \{\,dp \in \{1,\dots,G\}\mid G \bmod dp = 0\},
\]
\[
C' = T \times D = \{\, (tp,dp) \mid tp\in T,\; dp\in D\}.
\]
For each \((tp,dp)\in C'\), set
\[
G_{\mathrm{sim}} = \frac{G}{dp},
\qquad
ep = \Bigl\lceil \tfrac{n_r}{G_{\mathrm{sim}}}\Bigr\rceil,
\]
and define the simulator configuration 
\[
c = \bigl(G_{\mathrm{sim}},\,tp,\,ep\bigr)\in C.
\]

Then, for each \(c\in C\) and sequence length \(s\in S\), the simulator returns
\[
R_{\mathrm{prefill}}(s;\,c),
\]
and under \(dp\)-way data parallelism the effective total prefill throughput is
\[
R_{\mathrm{prefill}}(s;\,dp,c)
= dp \times R_{\mathrm{prefill}}(s;\,c)
= dp \times R_{\mathrm{prefill}}\!\Bigl(s;\,\tfrac{G}{dp},\,tp,\,ep\Bigr).
\]

However, before calculating the best parallelism configuration out of the set, we also need to decide which batch size to use.
Since we would like to optimise our configuration for overall throughput, we need to maximise the batch size for each input sequence length, such that we remain within the memory constraints of the island.
Let \(\max\_bs(s; c)\) denote the maximum batch size returned by the simulator for sequence length \(s\) under configuration \(c\).
We then round down to the nearest multiple of 8 and clamp between 8 and 1024:

\[
\mathrm{bs}(s;c)
= \min\!\Bigl(\max\bigl(\max\_bs(s;c) - \bigl(\max\_bs(s;c) \bmod 8\bigr),\,8\bigr),\,1024\Bigr).
\]

Similar to prefill, we seek a single “best” decode parallelism configuration across sequence lengths. 
For each simulator configuration \(c\in C\) and sequence length \(s\), we first compute the batch size \(\mathrm{bs}(s;c)\) as before, then measure the decode throughput $R_{\mathrm{decode}}\bigl(s, d;\,c, \mathrm{bs}(s;c)\bigr)$.
We define the optimal decode throughput at length \(s\) by
\[
R_{\mathrm{decode}}^*(s,d) = \max_{c\in C} R_{\mathrm{decode}}\bigl(s, d;\,c, \mathrm{bs}(s;c)\bigr),
\]
and the relative deviation of configuration \(c\) at \(s\) by
\[
\delta_{\mathrm{decode}}(c,s) = \frac{R_{\mathrm{decode}}^*(s,d) - R_{\mathrm{decode}}\bigl(s, d;\,c, \mathrm{bs}(s;c)\bigr)}{R_{\mathrm{decode}}^*(s,d)}.
\]
Aggregating these deviations over all lengths gives
\[
\begin{aligned}
\bar{\delta}_{\mathrm{decode}}(c) &= \frac{1}{\lvert S\rvert} \sum_{s\in S} \delta_{\mathrm{decode}}(c,s), 
&&\text{(uniform average over \(S\))},\\
\widetilde{\delta}_{\mathrm{decode}}(c) &= \sum_{s\in S} \delta_{\mathrm{decode}}(c,s)\,p(s\mid w),
&&\text{(weighted by workload distribution)}.
\end{aligned}
\]
We can then select the decode configuration that minimises the average deviation across all expected sequence lengths,
\[
c_{\mathrm{decode}} = \arg\min_{c\in C}\bar{\delta}_{\mathrm{decode}}(c).
\]

Shown in \autoref{fig:04-decode-throughput-vs-sequence-length}, for each GPU there is visibly a better parallelism configuration for most sequence lengths.
Unlike prefill, this configuration is not consistent across GPU types, but rather significantly impacted by the GPU specifications.
Configurations with zero throughput are due to the configuration exceeding the memory capabilities of the cluster.

\begin{figure}[ht]
    \centering
    \includegraphics[width=\textwidth]{simulator_decode_throughput_vs_seq_len.pdf}
    \caption[Stuff]
    {Stuff}
    \label{fig:04-decode-throughput-vs-sequence-length}
\end{figure}

The graphs of \autoref{fig:04-decode-throughput-vs-sequence-length} also demonstrate the advantages of heterogeneity.

% Add some explanations of the graph 

\subsection{Solver Resolution}
\label{chap:resolution}

We can use the simulator to benchmark the prefill and decode throughput of an island for a given input sequence length.
While it is feasible to consider each sequence length  (\(s\in S\)) individually, we can use a binning approach to reduce the search space of the scheduler.
The width of the ranges effectively control the resolution of the solution generated, which can be dynamically adjusted depending on input parameters.
This further reduces the solution space of the constraint optimisation problem.

We partition the sorted list of lengths \(S\) into equally‐sized intervals of width \(W\), called \textit{ranges}.
Rather than measure every \(s_i\), we sample a single representative length from each range at its midpoint.
With a positive integer width \(W\), we compute the midpoint offset as \(\mathrm{mid} = \left\lfloor \tfrac{W}{2}\right\rfloor\).
For each integer \(k\) such that \(0 \le k < \lfloor \lvert S\rvert / W\rfloor\), we let 
\[
\mathrm{range\_start}_k = k\,W,\quad
\mathrm{range\_end}_k = k\,W + W,\quad
\mathrm{range\_mid}_k = \mathrm{range\_start}_k + \mathrm{mid},
\]
and take \(s_{\mathrm{range\_mid}_k}\) as the representative length for the \(k\)-th range. We then compute the total probability of that range by summing 
\[
p_k = \sum_{i = \mathrm{range\_start}_k}^{\mathrm{range\_end}_k - 1} p\bigl(s_i \mid w\bigr),
\]
and record for range \(k\) the tuple \(\bigl(\mathrm{range\_mid}_k,\,p_k,\,\mathrm{range\_start}_k,\,\mathrm{range\_end}_k\bigr)\). Therefore, the set of all ranges is 
\[
\mathcal{R} = \bigl\{\,(\mathrm{range\_mid}_k,\,p_k,\,\mathrm{range\_start}_k,\,\mathrm{range\_end}_k)\mid 0 \le k < \lfloor \lvert S\rvert / W\rfloor \bigr\}.
\]
By sampling only at these midpoints and aggregating probabilities, we approximate the full set of sequence lengths in each range with a single point (\(\mathrm{range\_mid}_k\)), thereby greatly reducing the number of simulator calls while preserving workload‐weighted accuracy.

\subsubsection*{Prefill Throughput}

With the formed ranges \(\mathcal{R}\), each with their representative length \(s_r\) and total probability \(p_r\), we can redefine the allocation variables \(x_i(s)\) to act at the range level. 
Now island \(i\) assigns fraction \(x_i(r)\) of its GPU capacity to range \(r\in\mathcal{R}\), with
\[
\sum_{r\in\mathcal{R}} x_i(r) = 1,
\qquad
x_i(r)\ge0\quad\forall\,i\in I,\;r\in\mathcal{R}.
\]
Within range \(r\), every sequence length \(s\in S_r\) is approximated by \(s_r\), so that
\[
\sum_{s\in S_r} R_{\mathrm{prefill}}^{(i)}(s)\,p(s\mid w)\;\approx\;
R_{\mathrm{prefill}}^{(i)}(s_r)\;\sum_{s\in S_r}p(s\mid w)
\;=\;
R_{\mathrm{prefill}}^{(i)}(s_r)\,p_r.
\]
By including \(x_i(r)\) into our per‐range throughput models, the expected prefill throughput under workload \(w\) becomes
\[
\mathbb{E}[R_{\mathrm{prefill}}\mid w]
=\sum_{i\in I_{\mathrm{prefill}}}\sum_{r\in\mathcal{R}}
x_i(r)\;R_{\mathrm{prefill}}^{(i)}(s_r)\;p_r.
\]

\subsubsection*{Decode Throughput}

The same approach applies to the decode phase.
For each range \(r\) with representative length \(s_r\) and probability \(p_r\), we measure the decode throughput \(R_{\mathrm{decode}}^{(i)}\bigl(s_r,d\bigr)\) on island \(i\) for average decode length \(d\).
Island \(i\) allocates fraction \(x_i(r)\) of its capacity to \(r\), and within that range every sequence length is approximated by \(s_r\).  Thus,
\[
\sum_{s\in S_r} R_{\mathrm{decode}}^{(i)}(s,d)\,p(s\mid w)\;\approx\;
R_{\mathrm{decode}}^{(i)}(s_r,d)\;p_r.
\]
Normalising \(x_i(r)\) as before and summing over islands and ranges yields the expected decode throughput under workload \(w\) and decode length \(d\):
\[
\mathbb{E}[R_{\mathrm{decode}}\mid w,d]
=\sum_{i\in I_{\mathrm{decode}}}\sum_{r\in\mathcal{R}}
x_i(r)\;R_{\mathrm{decode}}^{(i)}\bigl(s_r,d\bigr)\;p_r.
\]
The overall expected throughput remains
\[
\mathbb{E}[R_{\mathrm{overall}}\mid w,d]
=\min\bigl(\mathbb{E}[R_{\mathrm{prefill}}\mid w],\;\mathbb{E}[R_{\mathrm{decode}}\mid w,d]\bigr).
\]

% Graph of binning techniques/resolution
%\begin{figure}[ht]
%    \centering
%    \includegraphics[width=\textwidth]{sample.pdf}
%    \caption[Stuff]
%    {Stuff}
%    \label{fig:04-sequence-length-binning}
%\end{figure}

%Single graph displaying throughput of prefill and decode for specific GPU (Not sure what this is for)

\clearpage
\section{Inner Solver}

The inner solver brings together the optimum parallelism section from \autoref{chap:prefill-parallelism} and \autoref{chap:decode-parallelism}.

Uses the constraint optimisation problem from \autoref{chap:constraint-optimisation}.


% STUFF
\clearpage
\section{Island Divider}

Why BO is difficult with high dimensionality, how we can reduce it with simple algorithm.

\clearpage
\subsection{Procedure}
The primary goal of the island divider procedure is to programmatically divide a given total number of GPUs of a specific type into a list of smaller, integer‐sized islands.
Let \(G_{\text{total}}\) denote the total number of GPUs available for that type, and let \(N_{\text{target}}\) represent the desired number of islands to create.
The actual number of islands may be lower if there are not enough GPUs to satisfy the minimum‐size requirement.
The parameter \(S_{\text{skew}}\) controls how any leftover GPUs (after each island has been assigned its minimum) are distributed: when \(S_{\text{skew}}\approx0.0\), the extra GPUs are spread as evenly as possible; if \(S_{\text{skew}}>0.0\), islands with higher indices receive proportionally more of the remainder; and if \(S_{\text{skew}}<0.0\), islands with lower indices receive more.
Finally, \(I_{\text{min}}\) specifies the minimum number of GPUs that any island must contain.

% we repeat this for every GPU type

\subsubsection{Preconditions}
Initially, we verify that the input parameters meet a series of constraints, if any of these conditions is not met, the procedure raises an error and terminates immediately.
The requirement \(I_{\text{min}} > 0\) ensures that each island contains at least one GPU, since a non‐positive minimum would be impossible.
Similarly, \(G_{\text{total}}\) must be nonnegative because a negative total GPU count is invalid, and \(N_{\text{target}}\) must also be nonnegative because a negative number of islands cannot be requested

Next, we perform trivial checks to assess whether it is possible to create at least one island given the input parameters.
If \(N_{\text{target}} = 0\), then no islands are to be created and the function returns the empty set \(\varnothing\).
The procedure then addresses cases in which allocation is not possible: if \(G_{\text{total}} = 0\), there are no GPUs to allocate and the function returns \(\varnothing\).
Similarly, if \(G_{\text{total}} < I_{\text{min}}\), there are insufficient GPUs to form even a single island of size \(I_{\text{min}}\), so the function again returns \(\varnothing\).
Only when all these preliminary checks have passed does the algorithm proceed to determine the feasible number of islands and distribute the GPUs accordingly.

% Maybe simplify to just say that bad inputs are discounted

\subsubsection{Number of Islands}

To establish how many islands can feasibly be created given the available GPUs and the minimum‐size requirement, we first compute \(N_{\text{max\_possible}} = \lfloor G_{\text{total}} / I_{\text{min}} \rfloor\), which represents the maximum number of islands if each island is allocated exactly \(I_{\text{min}}\) GPUs.
We then set \(N_{\text{actual}} = \min\bigl(N_{\text{target}},\,N_{\text{max\_possible}}\bigr)\), so that the actual number of islands cannot exceed either the user’s target or the physical limit imposed by the GPU pool.
If \(N_{\text{actual}} = 0\), this indicates that even a single island of size \(I_{\text{min}}\) cannot be formed (either because \(N_{\text{target}} = 0\) or because \(G_{\text{total}} < I_{\text{min}}\)), and the function returns the empty set \(\varnothing\).
Otherwise, \(N_{\text{actual}}\) islands will be created, each initially allocated \(I_{\text{min}}\) GPUs before any further distribution of the remaining GPUs occurs.

\subsubsection{Base Allocation}

In the base allocation step, we allocate \(I_{\text{min}}\) GPUs to each of the \(N_{\text{actual}}\) islands, so that the total GPUs used for this initial distribution is \(G_{\text{base\_used}} = N_{\text{actual}} \cdot I_{\text{min}}\).
The number of GPUs remaining to be distributed is then \(G_{\text{remaining}} = G_{\text{total}} - G_{\text{base\_used}}\).
At this point, every island has been provisioned with its minimum allocation of \(I_{\text{min}}\) GPUs, and the remaining \(G_{\text{remaining}}\) GPUs will be allocated in subsequent steps.

\subsubsection{Island Distribution}

In the case where \(G_{\text{remaining}} = 0\), all GPUs have already been assigned to the \(N_{\text{actual}}\) islands at their minimum size, so we make no weighted adjustments.
When \(N_{\text{actual}} = 1\), there is only a single island so we add all \(G_{\text{remaining}}\) GPUs to that island.
Otherwise we can distribute the remaining GPUs using the \(S_{\text{skew}}\)value.

The distribution begins with the calculation of a weight $w_k$ for each island $k$, where $k$ ranges from $1$ to $N_{\text{actual}}$, computed as $w_k = k^{S_{\text{skew}}}$.
However, if $S_{\text{skew}}$ is numerically close to zero, each island is given an equal weight of $w_k = 1.0$.
These weights are then normalised so that we can derive proportions $p_k$, indicating each island's share of the remaining GPUs. The sum of the weights is computed as
\[
W_{\text{sum}} = \sum_{j=1}^{N_{\text{actual}}} w_j.
\]
%If $W_{\text{sum}} \leq 0$, which may occur due to numerical errors or degenerate inputs, all islands are assigned equal proportions, $p_k = \frac{1}{N_{\text{actual}}}$. Otherwise, 
Each proportion is determined by dividing the weight of an island by the weight sum:
\[
p_k = \frac{w_k}{W_{\text{sum}}}.
\]
Using these proportions, we compute the ideal (possibly fractional) allocation $s_k^{\text{ideal}}$ for each island:
\[
s_k^{\text{ideal}} = p_k \cdot G_{\text{remaining}}.
\]

\subsubsection{Integer Adjustment}

Because GPUs can only be allocated in integer units, we use the largest remainder method to convert the fractional allocations into integers.
First, each island $k$ is provisionally assigned a base number of GPUs:
\[
s_k^{\text{base\_int}} = \left\lfloor s_k^{\text{ideal}} \right\rfloor.
\]
The current island sizes (initialised to $I_{\text{min}}$) are then increased by these integer allocations.

%\clearpage
%Next, the discrepancy between the total allocated GPUs and the intended $G_{\text{remaining}}$ is computed. This discrepancy,
%\[
%D = \text{round} \left( G_{\text{remaining}} - \sum_{j=1}^{N_{\text{actual}}} s_j^{\text{base\_int}} \right),
%\]
%represents the number of GPUs that are still unassigned due to rounding down the fractional values. To distribute the $D$ leftover GPUs, we calculate the fractional part for each island as
%\[
%f_k = s_k^{\text{ideal}} - s_k^{\text{base\_int}},
%\]
%and identify the $D$ islands with the largest fractional parts. Each of these selected islands receives one additional GPU.
%
%At the end of this process, each island has a final size given by
%\[
%S_k = I_{\text{min}} + s_k^{\text{base\_int}} + \delta_k,
%\]
%where $\delta_k \in \{0, 1\}$ depending on whether the island received an extra GPU during the discrepancy distribution. The total sum of all final island sizes equals $G_{\text{total}}$, satisfying the constraint that all available GPUs are used. The function returns this list of $N_{\text{actual}}$ island sizes as the result of the procedure.





% THIS IS USEFUL TEXT - BUT BETTER FOR LATER

%When profiling a specific GPU island for a given sequence length (for scheduling), the exact parallelism configuration of the GPUs must be known.
%Hence, we must first benchmark the GPU island across a range of sequence lengths that we expect to receive in our workload and determine the optimal parallelism configuration.
%While it is possible that the optimal parallelism configuration for a set of GPUs is different across a range of sequence lengths, we can assume that the best configurations are within a few percent of each other.
%We can integrate this benchmark parallelism selection step with the profiling step for the constraint solver to save computation requests

% Architecture
\section{System Architecture}
\label{chap:system-architecture}

%Furthermore, we statically define I to a set of curated islands, and the optimiser simply fits the workload to these islands.
%In order to generate the most optimal solution, we need to choose the best set of island sizes for each GPU type.
%However, adding this decision process to the constraint optimisation problem increases its complexity, as it effectively turns a bounded bin packing problem into an unbounded one.
%The optimiser must now determine not only the workload allocation but also the optimal sizes and number of islands.
%Hence, we need to formulate an architectural approach that incorporates this logic, which we detail in \autoref{chap:architecture}.

During \autoref{chap:constraint-optimisation} we formulated a non-linear bounded bin‐packing problem, by constraining the islands to a pre‐defined type and size.
In a true scheduling problem we expect to be allocated an inventory of different GPU types, let \(T\) be the set of GPU types.
For each \(t\in T\), let \(n_t\) be the number of GPUs of type \(t\) allocated in the inventory.
The full GPU inventory is then the vector \(\bigl(n_t\bigr)_{t\in T}\), with total number of GPUs in the system modelled as \(N = \sum_{t\in T} n_t\).

Similar to ThunderServe \autocite{jiangThunderServeHighperformanceCostefficient2025}, we can separate this scheduling problem as two separate, distinct stages.
The first stage is the island layout generation, which we can treat as a search‐like optimisation problem in which we explore combinations of GPU island sizes and types.
The second stage is the bounded bin‐packing step, a deterministic problem that assigns workload ranges to the pre‐computed islands to maximise throughput for a given workload.
By decoupling layout design from the packing phase, we can leverage specialised solvers for each subproblem and significantly reduce overall computational complexity. 

We refer to the bounded bin packing problem as the \textit{inner loop}, and the outer optimization process as the \textit{outer loop}, as illustrated in \autoref{fig:04-system-architecture}.
The \textit{outer loop} is responsible for generating promising island layouts, which are then passed to the \textit{inner loop} for resource packing.
The resulting throughput is returned to the \textit{outer loop}, allowing the optimisation process to evaluate and refine its search for effective configurations.
 This iterative cycle can be repeated multiple times to generate an optimal solution.
 
\begin{figure}[ht]
    \centering
   \begin{tikzpicture}[
			node distance=1.8cm and 2.5cm,
			every node/.style={draw, align=center, minimum height=1.2cm, minimum width=3cm},
			arrow/.style={->, thick}
		]
		
		% Nodes
		\node (outer) {Outer Optimiser};
		\node (inner) [right=of outer] {Inner min-cost\\ flow};
		\node (model) [below=of inner, yshift=-0.5cm] {Model};
	\end{tikzpicture}
    \caption[System Archiecture]
    {System Archiecture}
    \label{fig:04-system-architecture}
\end{figure}

\clearpage
\section{Inner Loop}

The \textit{inner loop} receives an island layout proposal from the \textit{outer loop}, this layout is specified as a finite index set \(\mathcal{I} \;=\; \{\,i = 1,\dots,|\mathcal{I}|\}\), where each island \(i\) is represented by a pair \(\bigl(t_{i},s_{i}\bigr)\).  
Here, \(t_{i}\in T\) denotes the GPU type assigned to island \(i\), and \(s_{i}\in \mathbb{Z}_{\ge 0}\) indicates the number of GPUs of type \(t_{i}\) allocated to that island.
The outer loop guarantees that these allocations respect the total‐inventory constraints.

Given the predefined island layout \(\bigl\{(t_{j},s_{j})\bigr\}_{j=1}^{|\mathcal{I}|}\), the inner loop’s only task is to assign workloads of different prefill and decode lengths.
No further modification of \(t_{i}\) or \(s_{i}\) takes place at this stage, as the inner loop focuses entirely on determining the optimal bin-packing configuration.

We can outline the steps of the \textit{inner loop} as follow: 

\begin{enumerate}
    \item \textbf{Form Ranges}  
    Take the specific workload distribution $w$, and separate into a series of ranges $\mathcal{R}$ using a specified resolution (\autoref{chap:resolution})
    
    \item \textbf{Prefill Configuration Generation \& Benchmark}  
    Identify the best prefill configuration ($c_{\mathrm{prefill}}$) for each island \(i\), and return the benchmark results describing the throughput for each range (\autoref{chap:prefill-parallelism}).
    
    \item \textbf{Decode Config Generation \& Benchmark}  
    Identify the best decode configuration ($c_{\mathrm{decode}}$) for each island \(i\), and return the benchmark results describing the throughput for each range (\autoref{chap:decode-parallelism}).
    
    \item \textbf{Problem Specification}  
    Formulate the exact optimisation problem based on the provided prefill and decode benchmarks (\autoref{chap:linear-problem}).
    
    \item \textbf{Solve}  
    Invoke the optimiser on the specified problem, to generate the best solution given the input parameters.
    
    \item \textbf{Results}  
    Extract the final workload‐assignment solution from the solver’s output, including any intermediary variables for debugging.
    
\end{enumerate}

% Many linear solution
\subsection{Linear Representation}

We can take the learning from the previous optimisation problem (\autoref{chap:constraint-optimisation}) to linearise our problem and combine prefill with the decode phases.
We can also take the learning from the solution resolution (\autoref{chap:resolution}), to reduce the number of variables and the problem complexity.


% THIS IS USEFUL TEXT - BUT BETTER FOR LATER - NOT SURE IF EVEN NEEDED

%When profiling a specific GPU island for a given sequence length (for scheduling), the exact parallelism configuration of the GPUs must be known.
%Hence, we must first benchmark the GPU island across a range of sequence lengths that we expect to receive in our workload and determine the optimal parallelism configuration.
%While it is possible that the optimal parallelism configuration for a set of GPUs is different across a range of sequence lengths, we can assume that the best configurations are within a few percent of each other.
%We can integrate this benchmark parallelism selection step with the profiling step for the constraint solver to save computation requests
% STUFF
\clearpage
\section{Outer Solver}

%We consider a resource allocation problem where the goal is to optimise a black-box objective function \( f: \mathcal{X} \to \mathbb{R} \) using Bayesian Optimisation.
%The function \( f \) is expensive to evaluate and only accessible via noisy observations. 
%The domain \( \mathcal{X} \subseteq \mathbb{R}^d \) represents a constrained space defined by hardware inventory limits.
%Our objective is to find an optimal configuration of GPU allocations that maximizes the function \( f \).

%
%\subsubsection*{Bayesian Optimization Formulation}
%
%Bayesian Optimization models the objective function \( f \) using a probabilistic surrogate model \( \hat{f}(x) \), typically a Gaussian Process (GP):
%
%\[
%\hat{f}(x) \sim \mathcal{GP}(\mu(x), k(x, x'))
%\]
%
%An acquisition function \( \alpha(x; \hat{f}) \) is used to determine the next point to evaluate:
%
%\[
%x_{t+1} = \arg \max_{x \in \mathcal{X}} \alpha(x; \hat{f}_t)
%\]
%
%This optimization is subject to a set of resource constraints, which are defined below.

%We refer to the bounded bin packing problem as the \textit{inner loop}, and the outer optimisation process as the \textit{outer loop}, as illustrated in \autoref{fig:04-problem-formulation}.
%The \textit{outer loop} is responsible for generating promising island layouts, which are then passed to the \textit{inner loop} for resource packing.
%The resulting throughput is returned to the \textit{outer loop}, allowing the optimisation process to evaluate and refine its search for effective configurations.
% This iterative cycle can be repeated multiple times to generate progressively generate an optimum solution.
 
% During \autoref{chap:constraint-optimisation} we formulated a non-linear bounded bin-packing problem, by constraining the islands to a pre-defined type and size.
%In a true scheduling problem we expect to be allocated an inventory of different GPU types, not formed islands.
%Let \(T\) be the set of GPU types an for each \(t\in T\), let \(n_t\) be the number of GPUs of type \(t\) allocated in the inventory.
%The full GPU inventory is then the vector \(\bigl(n_t\bigr)_{t\in T}\), with total number of GPUs in the system modelled as \(N = \sum_{t\in T} n_t\).

%The \textit{inner loop} receives an island layout proposal from the \textit{outer loop}, this layout is specified as a finite index set \(\mathcal{I} \;=\; \{\,i = 1,\dots,|\mathcal{I}|\}\), where each island \(i\) is represented by a pair \(\bigl(t_{i},s_{i}\bigr)\).  
%Here, \(t_{i}\in T\) denotes the GPU type assigned to island \(i\), and \(s_{i}\in \mathbb{Z}_{\ge 0}\) indicates the number of GPUs of type \(t_{i}\) allocated to that island.
%The \textit{outer loop} guarantees that these allocations respect the total-inventory constraints.

%Given the predefined island layout \(\bigl\{(t_{j},s_{j})\bigr\}_{j=1}^{|\mathcal{I}|}\), the inner loop's only task is to assign workloads of different prefill and decode lengths.
%No further modification of \(t_{i}\) or \(s_{i}\) takes place at this stage, as the inner loop focuses entirely on determining the optimal bin-packing configuration.

The \textit{outer loop} receives an input inventory, containing a list of GPU types (\(t\in T\)) and quantities (\(n_t\)), represented with the set \(\{\,I_t : t \in T\}\).
Given the inventory, the \textit{outer loop's} job is to propose new island layouts (\(\mathcal{I}\)), which can be subsequently evaluated by the \textit{inner loop}.

The \textit{inner loop} execution is somewhat expensive, hence a complete grid search would incur significant compute time.
Instead we use a refined Bayesian Optimisation search method with two different methods for searching, \textit{direct} (\autoref{chap:island-direct} and \textit{divided} (\autoref{chap:island-divide}).

\begin{figure}[ht]
    \centering
    \includegraphics[width=\textwidth,height=5cm]{example-image-a}
    \caption[Diagram showing effective purpose of \textit{outer loop}]
    {Diagram showing effective purpose of \textit{outer loop}.}
    \label{fig:04-outer-loop-demo}
\end{figure}

We select Bayesian Optimisation specifically due to its ability to efficiently explore high-dimensional, expensive‐to‐evaluate spaces with only a few hundred evaluations, rather than the exponential cost of a full grid search.
By maintaining a surrogate model of the island‐layout performance (Gaussian process over layouts), Bayesian Optimisation balances exploration (sampling layouts in poorly understood regions) with exploitation (refining around the current best layouts).

Unlike reinforcement learning, which often requires thousands of interactions to learn a policy, BO is designed for “one‐shot” black‐box optimisation.
With the correct parameters it can deliver near‐optimal solutions with far fewer simulator calls, minimal hyper-parameter tuning, and strong convergence.

% First BO Method
\subsection{Direct Island Generation}
\label{chap:island-direct}

In the \textit{direct} island generation approach, islands are explicitly defined within the Bayesian Optimisation process.
Configurations correspond to a fixed number of islands, with each candidate solution representing a specific allocation strategy.
The number of GPUs in each island is directly part of the output space, where specific variables represent the GPU count per island `slot'.
This setup provides precise control over how GPUs are grouped, but it also requires well-defined constraints to ensure that all generated configurations are valid.

\subsubsection*{Search Parameters}

\begin{itemize}
    \item \( T \): the set of GPU types, e.g., \( T = \{\text{A100}, \text{V100}, \ldots\} \)
    \item \( K \): the number of GPU slots per type
    \item \( x_{t,k} \): the integer decision variable representing the number of GPUs of type \( t \in T \) assigned to slot \( k \in \{1, \ldots, K\} \)
    \item \( I_t \): the total inventory of GPU type \( t \in T \)
\end{itemize}

%Each parameter \( x_{k,t} \) is bounded:
%
%\[
%x_{g,k} \in \{0, 1, \ldots, I_t\}, \quad \forall t \in T, \; \forall k \in \{1, \ldots, K\}
%\]

\subsubsection*{Search Constraints}

For each GPU type \( t \in T \), we enforce that the total number of GPUs allocated across all \( K \) slots does not exceed the inventory \( I_t \):
\[
\sum_{k=1}^{K} x_{g,k} \leq I_t, \quad \forall t \in T
\]

\subsubsection*{Implementation}

\todo{
\begin{itemize}
\item Ax Implementation.
\item Why SAASBO may be required (high dimensionality)
\item Why this approach is challenging.
\end{itemize}
}

% Though BO has been recognized as a powerful approaches for a diversity of problem domains, it is also known that traditional BO is inherently vulnerable to the ”curse-of- dimensionality” [6] when the search space has a large number of dimensions. wl-446

%These constraints define a feasible subspace \( \mathcal{X} \subseteq \mathbb{R}^d \), over which the Bayesian Optimization is performed.

%\subsubsection*{Search Space Definition}
%
%In implementation (e.g., using Ax or BoTorch), the search space \( \mathcal{X} \) is constructed by defining:
%\begin{itemize}
%    \item A list of bounded integer parameters \( \{x_{g,k}\} \)
%    \item A set of linear \texttt{SumConstraint} objects that enforce the inventory limits
%\end{itemize}
%
%These constraints are handled internally by the BO engine when sampling candidate configurations during optimization.


% Second BO method
\clearpage
\subsection{Divider Island Generation}
\label{chap:island-divide}

The \textit{divider} island generation method uses a more abstract representation that substantially reduces the dimensionality of the BO problem.
Explicit allocation at the slot level requires one parameter per slot per GPU type, resulting in a parameter space of size \( |G| \times K \).
This approach becomes quickly computationally expensive due to the high number of search parameters and the need for constraints to ensure valid configurations.

To address this, the divider-based approach reduces the search space to just two parameters per GPU type: the total number of islands and a scalar that controls the balance of GPU allocation across those islands.
 This formulation results in a significantly smaller parameter space of size \( |G| \times 2 \), enabling more efficient optimisation and faster convergence.

An advantage of this method is that it avoids the need for explicit resource constraints, which are otherwise necessary to enforce inventory limits and prevent invalid configurations.
By shifting from low-level allocation to higher-level distribution strategies, divider island generation maintains control over performance-influencing characteristics while offering a more scalable and interpretable optimisation process.

%The primary goal of the island divider procedure is to programmatically divide a given total number of GPUs of a specific type into a list of smaller, integer-sized islands.

\subsubsection{Input Variables}
Let \(G_{\text{total}}\) denote the total number of GPUs available for that type, and let \(N_{\text{target}}\) represent the desired number of islands to create.
We can allow the actual number of islands may be lower if there are not enough GPUs to satisfy the minimum-size requirement.
The parameter \(S_{\text{skew}}\) controls how any leftover GPUs (after each island has been assigned its minimum) are distributed: when \(S_{\text{skew}}\approx0.0\), the extra GPUs are spread as evenly as possible; if \(S_{\text{skew}}>0.0\), islands with higher indices receive proportionally more of the remainder; and if \(S_{\text{skew}}<0.0\), islands with lower indices receive more.
Finally, \(I_{\text{min}}\) specifies the minimum number of GPUs that any island must contain.

% we repeat this for every GPU type

\begin{figure}[ht]
    \centering
    \includegraphics[width=\textwidth]{divider_skew_effects_overlay_islands.pdf}
    \caption[Graph displaying the effect of changing the \(S_{\text{skew}}\) value]
    {Graph displaying the effect of changing the \(S_{\text{skew}}\) value where \(G_{\text{total}}=100\) and \(N_{\text{target}}=5\).}
    \label{fig:04-skew-changes}
\end{figure}

\subsubsection{Preconditions}
Initially, we verify that the input parameters meet a series of constraints, if any of these conditions is not met, the procedure raises an error and terminates immediately.
We require \(I_{\text{min}} > 0\), \(G_{\text{total}}\) and \(N_{\text{target}}\) must be nonnegative.
Next, we perform trivial checks to assess whether it is possible to create at least one island given the input parameters.
If \(N_{\text{target}} = 0\), \(G_{\text{total}} = 0\) or \(G_{\text{total}} < I_{\text{min}}\), then no islands are to be created and the function returns the empty set \(\varnothing\).
Similarly, if , there are insufficient GPUs to form even a single island of size \(I_{\text{min}}\), so the function again returns \(\varnothing\).

\subsubsection{Number of Islands}

To establish how many islands can feasibly be created given the available GPUs and the minimum-size requirement, we first compute \(N_{\text{max\_possible}} = \lfloor G_{\text{total}} / I_{\text{min}} \rfloor\), which represents the maximum number of islands if each island is allocated exactly \(I_{\text{min}}\) GPUs.
We then set \(N_{\text{actual}} = \min\bigl(N_{\text{target}},\,N_{\text{max\_possible}}\bigr)\), so that the actual number of islands cannot exceed either the user's target or the physical limit imposed by the GPU pool.
If \(N_{\text{actual}} = 0\), this indicates that even a single island of size \(I_{\text{min}}\) cannot be formed (either because \(N_{\text{target}} = 0\) or because \(G_{\text{total}} < I_{\text{min}}\)), and the function returns the empty set \(\varnothing\).
Otherwise, \(N_{\text{actual}}\) islands will be created, each initially allocated \(I_{\text{min}}\) GPUs before any further distribution of the remaining GPUs occurs.

\subsubsection{Base Allocation}

In the base allocation step, we allocate \(I_{\text{min}}\) GPUs to each of the \(N_{\text{actual}}\) islands, so that the total GPUs used for this initial distribution is \(G_{\text{base\_used}} = N_{\text{actual}} \cdot I_{\text{min}}\).
The number of GPUs remaining to be distributed is then \(G_{\text{remaining}} = G_{\text{total}} - G_{\text{base\_used}}\).
At this point, every island has been provisioned with its minimum allocation of \(I_{\text{min}}\) GPUs, and the remaining \(G_{\text{remaining}}\) GPUs will be allocated in subsequent steps.

\subsubsection{Island Distribution}

In the case where \(G_{\text{remaining}} = 0\), all GPUs have already been assigned to the \(N_{\text{actual}}\) islands at their minimum size, so we make no weighted adjustments.
When \(N_{\text{actual}} = 1\), there is only a single island so we add all \(G_{\text{remaining}}\) GPUs to that island.
Otherwise we can distribute the remaining GPUs using the \(S_{\text{skew}}\)value.

We first compute a weight $w_k$ for each island $k$, where $k$ ranges from $1$ to $N_{\text{actual}}$, computed as $w_k = k^{S_{\text{skew}}}$.
However, if $S_{\text{skew}}$ is numerically close to zero, each island is given an equal weight of $w_k = 1.0$.
These weights are then normalised so that we can derive proportions $p_k$, indicating each island's share of the remaining GPUs. The sum of the weights is computed as
\[
W_{\text{sum}} = \sum_{j=1}^{N_{\text{actual}}} w_j.
\]
%If $W_{\text{sum}} \leq 0$, which may occur due to numerical errors or degenerate inputs, all islands are assigned equal proportions, $p_k = \frac{1}{N_{\text{actual}}}$. Otherwise, 
Each proportion is determined by dividing the weight of an island by the weight sum:
\[
p_k = \frac{w_k}{W_{\text{sum}}}.
\]
Using these proportions, we compute the ideal (possibly fractional) allocation $s_k^{\text{ideal}}$ for each island:
\[
s_k^{\text{ideal}} = p_k \cdot G_{\text{remaining}}.
\]

\subsubsection{Integer Adjustment}

Because GPUs can only be allocated in integer units, we use the largest remainder method to convert the fractional allocations into integers.
First, each island $k$ is provisionally assigned a base number of GPUs:
\[
s_k^{\text{base\_int}} = \left\lfloor s_k^{\text{ideal}} \right\rfloor.
\]
The current island sizes (initialised to $I_{\text{min}}$) are then increased by these integer allocations.

%\todo{
%
%Next, the discrepancy between the total allocated GPUs and the intended $G_{\text{remaining}}$ is computed. This discrepancy,
%\[
%D = \text{round} \left( G_{\text{remaining}} - \sum_{j=1}^{N_{\text{actual}}} s_j^{\text{base\_int}} \right),
%\]
%represents the number of GPUs that are still unassigned due to rounding down the fractional values. To distribute the $D$ leftover GPUs, we calculate the fractional part for each island as
%\[
%f_k = s_k^{\text{ideal}} - s_k^{\text{base\_int}},
%\]
%and identify the $D$ islands with the largest fractional parts. Each of these selected islands receives one additional GPU.
%
%At the end of this process, each island has a final size given by
%\[
%S_k = I_{\text{min}} + s_k^{\text{base\_int}} + \delta_k,
%\]
%where $\delta_k \in \{0, 1\}$ depending on whether the island received an extra GPU during the discrepancy distribution. The total sum of all final island sizes equals $G_{\text{total}}$, satisfying the constraint that all available GPUs are used. The function returns this list of $N_{\text{actual}}$ island sizes as the result of the procedure.
%
%}


\clearpage
\subsubsection*{Search Parameters}

\begin{itemize}
    \item \( T \): the set of GPU types, e.g., \( T = \{\text{A100}, \text{V100}, \ldots\} \)
    \item \( n_t \): the target number of islands for GPU type \( t \in T \)
    \item \( s_t \): a vector value representing the skew or balance in GPU allocation across the \( n_t \) islands for GPU type \( t \)
    \item \( I_t \): the total available inventory of GPU type \( t \)
\end{itemize}

%Here, \( n_g \) is an integer representing how many islands to distribute GPUs across, and \( s_g \in [0, 1] \) controls how evenly the GPUs are distributed across those islands. For example, \( s_g = 0 \) could represent perfectly balanced islands (equal GPU count per island), while \( s_g = 1 \) could represent maximal skew (e.g., one island with most of the GPUs, others with minimal or none).

These two parameters (\( n_t \) and \( s_t \)) define the entire resource distribution strategy, allowing the Bayesian Optimisation process to search over high-level allocation patterns while implicitly satisfying inventory constraints by design.

\subsubsection*{Implementation}

\todo{
\begin{itemize}
\item Ax Implementation.
\item SAASBO not required.
\end{itemize}
}

\let\clearpage\origclearpage